% Einheit 1 von 3 – Zählprinzip und Anordnungen (bank/kombinatorik/e1.jsonl, zusammenbau v0.1)
%% TODO Zeitmarke Einheit 1: Marken-Zeile nicht lesbar
%% TODO Zweigzeile Teil 1: was hier gelernt wird (Satz in Schülersprache aus dem Einheitstitel) – Einheit 1
\einheitenkopf[e1]{Einheit 1 von 3 · Zählprinzip und Anordnungen}
\zweigzeile{Abitur LK}
\setcounter{aufgabe}{7}

%% TODO Titel: Ich-kann-Satz fehlt in der Bank (Platzhalter: Kettenname) – Nr. 8
%% TODO Anweisung: ein Satz über der ersten Teilaufgabe fehlt in der Bank – Nr. 8
\begin{aufgabe}{Zählprinzip und Anordnungen}
%% TODO \rechenplatz{n}: Schrittzahl des Grundfalls fehlt in der Bank (nur bei fester Schreibform, 2.3 b)
\begin{teile}
\teil Aus sieben Bewerbern werden Klassensprecher und Stellvertreter gewählt. Zählt die Reihenfolge? \\ \kreuz{ja} \\ \kreuz{nein}
\teil Ein Bäcker bietet vier Brotsorten und sieben Beläge an. Wie viele belegte Brote mit einer Sorte und einem Belag gibt es? \leerfeld
\teil Ein Restaurant bietet fünf Vorspeisen, acht Hauptgänge und einige Desserts an; so ergeben sich 160 dreigängige Menüs. Wie viele Desserts gibt es? \leerfeld
\teil Fünf verschiedene Bücher werden nebeneinander ins Regal gestellt. Wie viele Reihenfolgen gibt es? \leerfeld
\teil Wie viele verschiedene Buchstabenfolgen lassen sich aus den Buchstaben des Wortes ANANAS legen? \leerfeld
\teil Eine PIN besteht aus vier Ziffern von 0 bis 9, Wiederholung erlaubt. Wie viele PINs gibt es, und welcher Anteil davon besteht nur aus ungeraden Ziffern? \leerfeld PINs, Anteil \leerfeld
\teil Bei einem Lauf mit acht Teilnehmern werden Gold, Silber und Bronze vergeben. Wie viele Siegerlisten gibt es? \leerfeld
\teil Fünf verschiedenfarbige Buttons werden verteilt. Bestimme die Anzahl der Möglichkeiten, wenn I jedes von fünf Kindern genau einen Button bekommt, II jedes von fünf Kindern mehrere oder keinen bekommen kann, III jedes von acht Kindern höchstens einen bekommt \hfill (FHR 2020). I \leerfeld , II \leerfeld , III \leerfeld
\end{teile}
\end{aufgabe}

%% TODO Titel: Ich-kann-Satz fehlt in der Bank (Platzhalter: Kettenname) – Nr. 9
%% TODO Anweisung: ein Satz über der ersten Teilaufgabe fehlt in der Bank – Nr. 9
\begin{aufgabe}{Zählprinzip und Anordnungen – Fehler finden, Begründen, Anwendung}
\begin{teile}
\teil Drei verschiedene Aufkleber werden an drei Kinder verteilt; jedes Kind darf mehrere oder keinen bekommen. Mia rechnet so. Finde den Fehler und rechne richtig. \rechnung{3! &= 3 \cdot 2 \cdot 1 \\ &= 6}
\teil Begründe, warum bei einer Anordnung von lauter verschiedenen Dingen die Fakultät entsteht.
\teil Ein Fahrradschloss hat vier Ringe mit je sechs Ziffern. Ein Dieb probiert alle drei Sekunden eine Kombination. Wie lange braucht er höchstens, bis das Schloss offen ist? Gib die Zeit in Minuten an. \leerfeld[min]
\end{teile}
\end{aufgabe}
